% problem-000314 4 冠醚超分子自组装
\section*{4. 冠醚超分子自组装}
近期化学家合成了⼀种结构⾮常特殊的超分⼦，其设计思路是，利⽤分⼦⾃我识别与⾃我组装的能⼒，将3个环状化合物套⼊Y型分⼦的⻣架上，再通过化学反应将3个环状化合物的侧链进⾏关环形成第4个环。图⽰如下：

（图：）

4-1 已知化合物A、B的结构分别如下图所⽰，试问：它们通过什么作⽤进⾏⾃我识别与⾃我组装？
（图：）

\subsection*{4-2}
在⽤邻苯⼆酚和1,8-⼆氯-3,6-⼆氧杂⾟烷进⾏[2+2]环化合成化合物A中的⼤环冠醚部分时，需⽤到2种⽅法以提⾼成环产率，即⾼度稀释和加⼊⾦属离⼦（模板离⼦）。试说明为什么加⼊模板离⼦会提⾼产率？对Na+、K+、Rb+、Cs+⽽⾔哪种离⼦作模板产率最⾼？为什么？

\subsection*{4-3}
将超分⼦C转变为D时使⽤的化学反应通常称之为烯烃复分解反应，⽣成D的同时，还⽣成另⼀种化合物 $\mathbf { E } _ { \circ }$ 。已知D中的第4个环为三⼗三元环，试问，化合物E是什么？反应过程中每⽣成⼀个D同时⽣成⼏个E？

\subsection*{4-4}
K是合成化合物A的原料之⼀，以E为主要原料合成K的路线如下，请写出F、G、H、I、J、K所代表的化合物或反应条件。

$
\mathrm{E} \xrightarrow [ \mathrm{CCl} _ {4} ]{\mathrm{F}} \mathrm{G} \xrightarrow [ \mathrm{Br} _ {2} ]{\mathrm{H}} \frac {(\mathrm{CH} _ {3}) _ {3} \mathrm{CONa}}{(\mathrm{CH} _ {3}) _ {3} \mathrm{COH}} \xrightarrow [ \mathrm{I} ]{\mathrm{Mg}} \mathrm{J} \xrightarrow [ \mathrm{THF} ]{1 . \mathrm{G}} \mathrm{K}
$

\subsection*{4-5}
K也可以⽤如下所⽰的路线合成，写出 $\mathbf { L } _ { \mathbf { \nu } }$ 、M、N、O所代表的化合物。

$
\mathrm{L} \xrightarrow {\mathrm{H} _ {2} \mathrm{O}} \mathrm{M} \xrightarrow {\mathrm{NaNH} _ {2}} \mathrm{N} \xrightarrow [ 2 . \mathrm{H} _ {2} \mathrm{O} ]{1 . \mathrm{G}} \mathrm{O} \xrightarrow [ \text {Lindlar Pd} ]{\mathrm{H} _ {2}} \mathrm{K}
$
